\documentclass[11pt]{article}
\usepackage[localtheorems]{raeez-math-template}

\newtheorem{theorem}{Theorem}
\newtheorem{proposition}[theorem]{Proposition}
\newtheorem{lemma}[theorem]{Lemma}
\newtheorem{corollary}[theorem]{Corollary}
\newtheorem{conjecture}[theorem]{Conjecture}
\theoremstyle{definition}
\newtheorem{definition}[theorem]{Definition}
\newtheorem{remark}[theorem]{Remark}

\providecommand{\C}{\mathbb{C}}
\providecommand{\Z}{\mathbb{Z}}
\providecommand{\Q}{\mathbb{Q}}
\providecommand{\cO}{\mathcal{O}}
\providecommand{\cM}{\mathcal{M}}
\providecommand{\cF}{\mathcal{F}}
\providecommand{\cZ}{\mathcal{Z}}
\providecommand{\cH}{\mathcal{H}}
\providecommand{\cC}{\mathcal{C}}
\providecommand{\cT}{\mathcal{T}}
\providecommand{\fg}{\mathfrak{g}}
\providecommand{\fgl}{\mathfrak{gl}}
\providecommand{\Perf}{\mathrm{Perf}}
\providecommand{\Coh}{\mathrm{Coh}}
\providecommand{\End}{\mathrm{End}}
\providecommand{\Hilb}{\mathrm{Hilb}}
\providecommand{\Stab}{\mathrm{Stab}}
\providecommand{\Obs}{\mathrm{Obs}}
\providecommand{\hCS}{\mathrm{hCS}}
\providecommand{\CoHA}{\mathrm{CoHA}}
\providecommand{\Rep}{\mathrm{Rep}}
\providecommand{\DT}{\mathrm{DT}}
\providecommand{\GW}{\mathrm{GW}}
\providecommand{\PT}{\mathrm{PT}}
\providecommand{\GV}{\mathrm{GV}}
\providecommand{\BPS}{\mathrm{BPS}}
\providecommand{\Tr}{\mathrm{Tr}}
\providecommand{\ch}{\mathrm{ch}}
\providecommand{\Li}{\mathrm{Li}}
\providecommand{\MM}{\mathrm{M}}
\providecommand{\vir}{\mathrm{vir}}
\providecommand{\red}{\mathrm{red}}
\providecommand{\topc}{\mathrm{top}}
\providecommand{\Eone}{E_1}
\providecommand{\Etwo}{E_2}
\providecommand{\Ethree}{E_3}
\providecommand{\PhiFA}{\Phi^{\mathrm{FA}}}
\providecommand{\SpCh}{\mathrm{Sp}^{\mathrm{ch}}}
\providecommand{\kcat}{\kappa_{\mathrm{cat}}}
\providecommand{\kch}{\kappa_{\mathrm{ch}}}
\providecommand{\kchcomp}{\kappa_{\mathrm{ch}}^{\mathrm{comp}}}
\providecommand{\kchHeis}{\kappa_{\mathrm{ch}}^{\mathrm{Heis}}}
\providecommand{\kBKM}{\kappa_{\mathrm{BKM}}}
\providecommand{\kfib}{\kappa_{\mathrm{fiber}}}
\providecommand{\DeltaFive}{\Delta_5}
\providecommand{\PhiTen}{\Phi_{10}}
\providecommand{\Dfive}{D_5}
\providecommand{\Wilpha}{\mathcal{W}_{1+\infty}}

\title{MNOP, Wall Crossing, and the Trace Comparison}
\author{Raeez Lorgat}
\date{2026-04-30}

\begin{document}
\maketitle

\begin{abstract}
The Maulik-Nekrasov-Okounkov-Pandharipande comparison relates three
different curve-counting theories on a Calabi-Yau threefold: ideal-sheaf
Donaldson-Thomas theory, stable-pair theory, and Gromov-Witten stable-map
theory.  Its equality has content only after the degree-zero DT factor is
removed and the variables are identified by
\[
  -q=e^{iu}.
\]
Kontsevich-Soibelman wall crossing supplies the DT/PT part of the
comparison.  The PT/GW part is the BPS or Gopakumar-Vafa
reorganisation of stable-pair and stable-map series.  A trace
formulation packages the three numerical series as character images of
traces in a common \(\Etwo\)-centre only after dualisable modules,
orientation data, quantum-master-equation completions, and comparison
maps have been supplied.

For \(K3\times E\), the reduced primitive scalar uses the
Oberdieck-Pandharipande normalisation
\[
  \Dfive=64^{-1}\DeltaFive,\qquad
  \PhiTen^{\mathrm{OP}}=\Dfive^2=4096^{-1}\DeltaFive^2,\qquad
  Z^{\DT,\red,\prime}_{K3\times E}
  =-\bigl(\PhiTen^{\mathrm{OP}}\bigr)^{-1}
  =-4096\,\DeltaFive^{-2}.
\]
This scalar is separate from
\[
  \kcat(K3\times E)=0,\quad
  \kchcomp(K3\times E)=0,\quad
  \kchHeis(K3\times E)=3,\quad
  \kBKM(\fg_{\DeltaFive})=5,\quad
  \kfib(K3)=24.
\]
\end{abstract}

\tableofcontents

\section{The Three Curve-Counting Series}

Let \(X\) be a smooth Calabi-Yau threefold and
\(\beta\in H_2(X,\Z)\) an effective curve class.  Compact projective
targets, toric local targets, and reduced \(K3\times E\) targets use
different properness and equivariance packages; the moduli problems
remain distinct.

\begin{definition}[Donaldson-Thomas series]
The ideal-sheaf moduli space
\[
  I_n(X,\beta)=\Hilb(X,\beta,n)
\]
parametrises ideal sheaves \(I_Z\subset \cO_X\) with
\([Z]=\beta\) and \(\chi(\cO_Z)=n\).  In the proper or equivariantly
proper setting, the symmetric obstruction theory and Behrend function
\(\nu_I\) define
\[
  \DT_{n,\beta}(X)
  =
  \int_{[I_n(X,\beta)]^{\vir}}1
  =
  \chi\bigl(I_n(X,\beta),\nu_I\bigr).
\]
The partition function and its degree-zero reduction are
\[
  Z_X^{\DT}(q,Q)
  =
  \sum_{\beta,n}\DT_{n,\beta}(X)\,q^nQ^\beta,
  \qquad
  Z_X^{\DT,\prime}(q,Q)
  =
  \frac{Z_X^{\DT}(q,Q)}{Z_{X,0}^{\DT}(q)}.
\]
For a projective Calabi-Yau threefold,
\[
  Z_{X,0}^{\DT}(q)=\MM(-q)^{\chi_{\topc}(X)},\qquad
  \MM(q)=\prod_{m\ge1}(1-q^m)^{-m}.
\]
\end{definition}

\begin{definition}[Pandharipande-Thomas series]
The stable-pair moduli space \(P_n(X,\beta)\) parametrises pairs
\[
  \cO_X\longrightarrow F
\]
with \(F\) pure of dimension one, \([F]=\beta\), and \(\chi(F)=n\).
Its virtual class defines
\[
  \PT_{n,\beta}(X)=\int_{[P_n(X,\beta)]^{\vir}}1,
  \qquad
  Z_X^{\PT}(q,Q)=
  \sum_{\beta,n}\PT_{n,\beta}(X)\,q^nQ^\beta .
\]
Joyce-Song, Toda, and Bridgeland wall crossing give
\[
  Z_X^{\DT,\prime}(q,Q)=Z_X^{\PT}(q,Q)
\]
in their stated Calabi-Yau threefold settings and chambers.
\end{definition}

\begin{definition}[Gromov-Witten series]
The stable-map stack \(\overline{\cM}_{g,0}(X,\beta)\) carries a virtual
class of expected dimension zero.  The primary invariant is
\[
  \GW_{g,\beta}(X)
  =
  \int_{[\overline{\cM}_{g,0}(X,\beta)]^{\vir}}1 .
\]
After separating constant maps, the disconnected series is
\[
  Z_X^{\GW,\prime}(u,Q)
  =
  \exp\!\left(
    \sum_{\beta>0}\sum_{g\ge0}
      \GW_{g,\beta}(X)\,u^{2g-2}Q^\beta
  \right).
\]
\end{definition}

\begin{theorem}[MNOP comparison]
\label{thm:mnop-comparison}
On theorem loci of the DT/GW correspondence,
\[
  Z_X^{\GW,\prime}(u,Q)
  =
  Z_X^{\DT,\prime}(q,Q)\big|_{-q=e^{iu}}
  =
  Z_X^{\PT}(q,Q)\big|_{-q=e^{iu}}.
\]
The theorem loci include the smooth toric threefolds in the
topological-vertex setting of MNOP and MOOP, toric local examples such as
the resolved conifold, and the reduced primitive \(K3\times E\) theory
of Oberdieck-Pandharipande, Oberdieck-Pixton, and Oberdieck.  For a
generic smooth projective Calabi-Yau threefold this displayed equality
is the MNOP conjecture.
\end{theorem}

\begin{remark}[The variable]
The substitution \(-q=e^{iu}\) is part of the theorem data.  The sign
belongs to the Behrend-weighted DT convention.  The exponential records
the passage from connected stable-map free energies to the disconnected
GW partition function.
\end{remark}

\section{Kontsevich-Soibelman Wall Crossing}

Let \(\cC=\Perf(X)\) and
\(\Gamma=K^{\mathrm{num}}_0(\cC)\), equipped with the skew Euler form
\[
  \langle \gamma_1,\gamma_2\rangle
  =
  \chi_\cC(\gamma_1,\gamma_2)-\chi_\cC(\gamma_2,\gamma_1).
\]
The completed quantum torus has generators \(\hat e_\gamma\) and product
\[
  \hat e_{\gamma_1}\hat e_{\gamma_2}
  =
  q^{\langle\gamma_1,\gamma_2\rangle/2}
  \hat e_{\gamma_1+\gamma_2}.
\]
For a primitive class \(\gamma\), the quantum dilogarithm operator is
\[
  \Psi_\gamma^q
  =
  \prod_{n\ge0}\bigl(1-q^{n+1/2}\hat e_\gamma\bigr)^{-1},
\]
acting in the pro-nilpotent completion fixed by a strict sector in the
central-charge plane.  Its semiclassical logarithm is
\[
  \log \Psi_\gamma^q
  \sim
  \sum_{k\ge1}\frac{e_{k\gamma}}{k^2}
  =
  \Li_2(e_\gamma),
\]
with the usual KS sign twist in the \(q\to1\) limit.

\begin{theorem}[Kontsevich-Soibelman wall crossing]
\label{thm:ks-wall-crossing}
Let \(\sigma,\sigma'\in\Stab(\cC)\) be connected by a path crossing
walls \(W_{\gamma_1},\ldots,W_{\gamma_r}\) in a strict sector.  The
path-ordered product
\[
  U^\rho_{\sigma\to\sigma'}
  =
  \prod_{j=1}^r
  \left(\Psi^q_{\gamma_j}\right)^{\Omega_{\sigma_j}(\gamma_j)}
\]
is invariant under homotopy of the path with endpoints and sector fixed.
Equivalently, the ordered product of BPS automorphisms attached to a
sector is constant as the stability condition crosses its walls.
\end{theorem}

\begin{proof}
This is Kontsevich-Soibelman's motivic wall-crossing theorem.  The proof
uses the motivic Hall algebra, the integration map to the quantum torus,
and the quantum dilogarithm identities attached to rank-two wall
patterns.
\end{proof}

\begin{proposition}[DT/PT/GW factorisation]
\label{prop:dt-pt-gw-factorisation}
On every locus where the required wall-crossing and BPS expansions are
proved, the comparison in Theorem~\ref{thm:mnop-comparison} factors as
\[
  Z_X^{\DT}(q,Q)
  =
  \MM(-q)^{\chi_{\topc}(X)}Z_X^{\PT}(q,Q),
\]
followed by the BPS expansion of the PT series and the GV stable-map
summation
\[
  F_X^{\GW}(u,Q)
  =
  \sum_{\beta>0}\sum_{g\ge0}\sum_{k\ge1}
  n^g_\beta(X)\,
  \frac{1}{k}
  \bigl(2\sin\tfrac{ku}{2}\bigr)^{2g-2}
  Q^{k\beta}.
\]
For a fixed curve class, the PT BPS expansion is written in the
Pandharipande-Thomas convention as
\[
  \sum_n P_{n,\beta}q^n
  =
  \sum_{g\ge0} n^g_\beta\,q^{1-g}(1+q)^{2g-2}.
\]
Under \(-q=e^{iu}\), the factor \(q^{1-g}(1+q)^{2g-2}\) becomes the
corresponding sine factor up to the standard sign convention absorbed in
the BPS integer \(n^g_\beta\).
\end{proposition}

\begin{proof}
The first equality is the DT/PT wall-crossing theorem.  The displayed
BPS expansion is the PT rationality and integrality package of
Pandharipande-Thomas, with its later projective proofs in the available
generality.  The GV summation is the Gopakumar-Vafa expression of the
stable-map free energy in terms of the same BPS integers.  The
substitution \(-q=e^{iu}\) identifies the PT basis
\(q^{1-g}(1+q)^{2g-2}\) with the sine basis of the GW expansion.
\end{proof}

\section{Trace Formulation}

The two-stage CY-to-chiral construction has a factorisation-algebra
stage
\[
  \PhiFA_3(\Perf(X))=\cF_X.
\]
In the holomorphic field-theoretic lane one writes
\[
  \cF_X\simeq\Obs^q_{\hCS}(X)
\]
after solving the quantum master equation on the chosen completion of
fields.  The specialised CY-to-chiral output at dimension \(d\ge3\) is
native \(\Eone\).  The \(\Etwo\)-structure used by trace comparisons
lives on a centre, double, or module category, for example an
\(\Etwo\)-centre of \(\cF_X\)-modules.

\begin{definition}[\(\Ethree\)-trace target]
Let \(A\) be an \(\Ethree\)-algebra in a symmetric monoidal
\((\infty,1)\)-category and let \(M\) be a dualisable \(A\)-module.
Dunn-Lurie additivity gives a trace value
\[
  \Tr_A^{\Ethree}(M,\mathrm{id}_M)
  \in
  \End_A^{\Etwo}(\mathbf 1_A).
\]
A character map from this \(\Etwo\)-centre to a completed
generating-function ring turns the trace value into a numerical
partition function.
\end{definition}

\begin{theorem}[Conditional trace comparison]
\label{thm:conditional-trace-comparison}
Let \(X\) be a smooth compact Calabi-Yau threefold.  Assume:
\begin{enumerate}
\item the hCS or TCFT construction of \(\cF_X\) satisfies the quantum
master equation on the completed field space;
\item the DT, PT, and GW geometries define dualisable \(\cF_X\)-modules
\[
  M_\DT,\qquad M_\PT,\qquad M_\GW;
\]
\item their orientation data agree with the \((-1)\)-shifted symplectic
structures and with Behrend weighting on the DT side;
\item comparison maps identify trace insertions with ideal-sheaf,
stable-pair, and stable-map insertions;
\item the MNOP or PT/GW correspondence is established in the relevant
chamber.
\end{enumerate}
Then the character images of the three trace values satisfy
\[
  \ch\Tr_{\cF_X}^{\Ethree}(M_\DT)
  =
  \MM(-q)^{\chi_{\topc}(X)}
  \ch\Tr_{\cF_X}^{\Ethree}(M_\PT),
\]
and
\[
  \ch\Tr_{\cF_X}^{\Ethree}(M_\GW)
  =
  \left(\ch\Tr_{\cF_X}^{\Ethree}(M_\PT)\right)\big|_{-q=e^{iu}}.
\]
Thus the trace formulation is equivalent to MNOP after applying the
specified character maps.  Equality of full trace maps on all
observables is the stronger TCFT comparison.
\end{theorem}

\begin{proof}
The first displayed equality is DT/PT wall crossing in trace character
form, including the degree-zero DT factor.  The second is the PT/GW
comparison after the MNOP variable substitution.  The hypotheses place
the three character values in one completed \(\Etwo\)-centre and make
the enumerative integrals into the character images of the corresponding
trace insertions.
\end{proof}

\section{Affine and Local Tests}

\begin{proposition}[\(\C^3\)]
\label{prop:c3}
For \(X=\C^3\) with the standard Calabi-Yau torus action,
\[
  Z_{\C^3}^{\DT}(q)=\MM(-q).
\]
The equivariant hCS trace on the vacuum insertion has the same value:
\[
  \Tr^{\cF_{\C^3},T}_{\DT}(\mathbf1;q)
  =
  \sum_{\pi}(-q)^{|\pi|}
  =
  \MM(-q),
\]
where \(\pi\) ranges over plane partitions.  The GW vertex gives
\[
  Z_{\C^3}^{\GW}(u)=\MM(e^{iu}),
\]
and the two vacuum amplitudes agree under \(-q=e^{iu}\).
\end{proposition}

\begin{proof}
The fixed points of \(\Hilb^n(\C^3)\) are plane partitions.  The
Calabi-Yau specialisation of the virtual tangent ratio gives the Behrend
sign, hence the MacMahon series \(\MM(-q)\).  The topological vertex
computes the corresponding GW vacuum amplitude.
\end{proof}

The affine toric Hall algebra identity used here is
\[
  \CoHA(\C^3)\simeq Y^+(\widehat{\mathfrak{gl}}_1).
\]
The full \(\Wilpha\) algebra appears after the negative half, Cartan
factor, Hopf pairing, Drinfeld double or centre, completion, and
evaluation data are added.

\begin{proposition}[Resolved conifold]
\label{prop:conifold}
Let
\[
  X=\cO_{\mathbb P^1}(-1)\oplus\cO_{\mathbb P^1}(-1).
\]
The toric theorem gives
\[
  Z_{\mathrm{con}}^{\DT}(q,Q)
  =
  \MM(-q)^2
  \prod_{k\ge1}\bigl(1-Q(-q)^k\bigr)^k,
\]
so
\[
  Z_{\mathrm{con}}^{\PT}(q,Q)
  =
  \prod_{k\ge1}\bigl(1-Q(-q)^k\bigr)^k.
\]
The single nonzero BPS invariant is \(n^0_{[\mathbb P^1]}=1\), and the
GW free energy is
\[
  F_{\mathrm{con}}^{\GW}(u,Q)
  =
  \sum_{k\ge1}
  \frac{Q^k}{k(2\sin(ku/2))^2}.
\]
These formulas give MNOP for the conifold under \(-q=e^{iu}\).
\end{proposition}

\begin{proof}
The DT formula is the topological-vertex calculation, with
\(\chi_{\topc}(X)=2\).  Removing \(\MM(-q)^2\) gives the stable-pair
series.  The Bryan-Pandharipande GW formula is the GV summation with
one genus-zero BPS state in the compact curve class.
\end{proof}

\section{\texorpdfstring{\(K3\times E\)}{K3 x E}}

Let \(S\) be a K3 surface, \(E\) an elliptic curve, and \(X=S\times E\).
The translation action of \(E\) forces reduced obstruction theories.
For primitive nonzero \(\beta_h\in H_2(S,\Z)\) with
\(\langle\beta_h,\beta_h\rangle=2h-2\), the reduced DT, PT, and GW
series are compared in the Oberdieck-Pandharipande and
Oberdieck-Pixton theory.

\begin{proposition}[Reduced primitive scalar]
\label{prop:k3e-op-scalar}
In the theta normalisation of the Igusa square root,
\[
  [q^{1/2}r^{1/2}s^{1/2}]\DeltaFive=64,
  \qquad
  \Dfive=64^{-1}\DeltaFive.
\]
The OP Igusa form is
\[
  \PhiTen^{\mathrm{OP}}
  =
  \Dfive^2
  =
  4096^{-1}\DeltaFive^2,
\]
while the unnormalised Igusa square is
\[
  \PhiTen^{\mathrm{un}}=\DeltaFive^2.
\]
The primitive reduced branch is
\[
  Z_X^{\DT,\red,\prime}
  =
  Z_X^{\PT,\red}
  =
  Z_X^{\mathrm{OP}}
  =
  -\bigl(\PhiTen^{\mathrm{OP}}\bigr)^{-1}
  =
  -\Dfive^{-2}
  =
  -4096\,\DeltaFive^{-2}.
\]
\end{proposition}

\begin{proof}
Igusa's theta product gives
\([q^{1/2}r^{1/2}s^{1/2}]\DeltaFive=64\), hence
\(\Dfive=64^{-1}\DeltaFive\).  Oberdieck-Pixton use the monic square
\(\PhiTen^{\mathrm{OP}}=\Dfive^2\).  Oberdieck's reduced DT/PT theorem
identifies the reduced DT and PT quotient series, and the
Oberdieck-Pandharipande primitive comparison identifies this series with
the OP reduced GW expansion after \(P=e^{iu}\) and the DT sign convention
\((-P)^n\).  The displayed sign is the OP scalar convention.
\end{proof}

\begin{proposition}[Five separate \(K3\times E\) numbers]
\label{prop:k3e-five-numbers}
The following quantities have independent definitions:
\[
\begin{array}{c|c|c}
\text{quantity} & \text{value} & \text{source}\\
\hline
\kcat(K3\times E) & 0
& \chi(\cO_{K3})\chi(\cO_E)=2\cdot0\\
\kchcomp(K3\times E) & 0
& \sum_q(-1)^qh^{0,q}(K3\times E)=1-1+1-1\\
\kchHeis(K3\times E) & 3
& \text{Mukai-Heisenberg specialisation}\\
\kBKM(\fg_{\DeltaFive}) & 5
& c_1(0)/2=10/2\\
\kfib(K3) & 24
& \mathrm{rk}\,\widetilde H(K3,\Z)
\end{array}
\]
The scalar \(Z_X^{\mathrm{OP}}\) is partition-function data.  These five
numbers live in distinct constructions, and none is obtained by summing
the others.
\end{proposition}

\begin{proof}
The first row is Kunneth multiplicativity for \(\chi(\cO)\) on the total
space.  The second is the compact Hodge supertrace from the
\((1,1,1,1)\) Hodge column.  The third is a chiral Heisenberg-Mukai
specialisation.  The fourth is Borcherds' weight formula applied to the
Gritsenko form \(\DeltaFive\), whose Jacobi input has constant term
\(c_1(0)=10\).  The fifth is the Mukai-lattice rank of the K3 fibre.
\end{proof}

\section{Final Form}

\begin{theorem}[MNOP as wall crossing plus trace comparison]
\label{thm:final-form}
Let \(X\) be a Calabi-Yau threefold in one of the proved MNOP settings:
smooth toric threefolds in the vertex formalism, toric local examples
such as the conifold, or the reduced primitive \(K3\times E\) theory.
Then the numerical comparison is
\[
  Z_X^{\GW,\prime}(u,Q)
  =
  \left(
  \frac{Z_X^{\DT}(q,Q)}{\MM(-q)^{\chi_{\topc}(X)}}
  \right)\bigg|_{-q=e^{iu}}
  =
  Z_X^{\PT}(q,Q)\big|_{-q=e^{iu}},
\]
with reduced obstruction theories and OP variables in the \(K3\times E\)
case.  Kontsevich-Soibelman wall crossing accounts for the DT/PT
equality.  The PT/GW equality is the stable-pairs/GV/stable-map
comparison.  The trace reading of this identity is valid under the
dualisability, orientation, QME/completion, and comparison-map
hypotheses of Theorem~\ref{thm:conditional-trace-comparison}.
\end{theorem}

\begin{proof}
The first equality is the combination of degree-zero DT removal and
DT/PT wall crossing.  The second is the MNOP variable substitution and
the stable-pairs/GW comparison on the named theorem loci.  The trace
form follows by applying the character map to the three trace values in
the common \(\Etwo\)-centre supplied by the hypotheses.
\end{proof}

\section{References}

\begin{itemize}[leftmargin=1.1em]
\item D. Maulik, N. Nekrasov, A. Okounkov, and R. Pandharipande,
\emph{Gromov-Witten theory and Donaldson-Thomas theory} I-II,
Compositio Math. 142 (2006), arXiv:math/0312059 and arXiv:math/0406092.
\item D. Maulik, A. Oblomkov, A. Okounkov, and R. Pandharipande,
\emph{Gromov-Witten/Donaldson-Thomas correspondence for toric
3-folds}, Invent. Math. 186 (2011), arXiv:0809.3976.
\item R. Pandharipande and R. P. Thomas, \emph{Curve counting via stable
pairs in the derived category}, Invent. Math. 178 (2009),
arXiv:0707.2348, and the stable-pairs rationality papers.
\item T. Bridgeland, \emph{Hall algebras and curve-counting invariants},
J. Amer. Math. Soc. 24 (2011), arXiv:1002.4372.
\item Y. Toda, \emph{Curve counting theories via stable objects I},
J. Amer. Math. Soc. 23 (2010), arXiv:0902.4371.
\item M. Kontsevich and Y. Soibelman, \emph{Stability structures,
motivic Donaldson-Thomas invariants and cluster transformations},
arXiv:0811.2435.
\item R. Gopakumar and C. Vafa, \emph{M-theory and topological strings}
I-II, arXiv:hep-th/9809187 and arXiv:hep-th/9812127.
\item S. Katz, A. Klemm, and C. Vafa, \emph{M-theory, topological
strings and spinning black holes}, Adv. Theor. Math. Phys. 3 (1999).
\item J. Bryan and R. Pandharipande, \emph{The local Gromov-Witten
theory of curves}, arXiv:math/0411037.
\item G. Oberdieck and R. Pandharipande, \emph{Curve counting on
\(K3\times E\), the Igusa cusp form \(\chi_{10}\), and descendent
integration}, arXiv:1411.1514.
\item G. Oberdieck and A. Pixton, \emph{Holomorphic anomaly equations and
the Igusa cusp form conjecture}, Invent. Math. 213 (2018),
arXiv:1706.10100.
\item G. Oberdieck, \emph{On reduced stable pair invariants},
Duke Math. J. 167 (2018), arXiv:1510.06561.
\item R. E. Borcherds, \emph{Automorphic forms with singularities on
Grassmannians}, Invent. Math. 132 (1998).
\item V. Gritsenko and V. Nikulin, \emph{Siegel automorphic form
corrections of some Lorentzian Kac-Moody Lie algebras}, Amer. J. Math.
119 (1997), and the Gritsenko weight-five paramodular form.
\item K. Costello and S. Li, \emph{Quantum BCOV theory on Calabi-Yau
manifolds and the higher genus B-model}, arXiv:1201.4501, and
Costello-Gwilliam, \emph{Factorization algebras in quantum field theory}.
\item O. Schiffmann and E. Vasserot, the cohomological Hall algebra of
\(\C^3\) and \(Y^+(\widehat{\mathfrak{gl}}_1)\).
\end{itemize}

\end{document}
